\documentclass[12pt,a4paper,oneside]{article}
\usepackage[brazil]{babel}
\usepackage[utf8]{inputenc}
\usepackage{olymp}


\newlength{\exmpwidouf}
\exmpwidinf=10cm
\exmpwidouf=6cm


\setcounter{problem}{0}

\begin{document}

\begin{problem}{Idade da Camila}{idade}{2}

Cibele, Camila e Celeste são três irmãs inseparáveis. Estão sempre juntas e adoram fazer esportes, ler, cozinhar, jogar no computador... Agora estão aprendendo a programar computadores para desenvolverem seus próprios jogos.

Mas nada disso interessa para esta tarefa: estamos interessados apenas nas suas idades. Sabemos que Cibele nasceu antes de Camila e Celeste nasceu depois de Camila.

Dados três números inteiros indicando as idades das irmãs, escreva um programa para determinar a idade de Camila.

%\Notes
%
%Observações, se tiver.

\InputFile

A entrada é composta por três linhas, cada linha contendo um número inteiro $N$, a idade de uma das irmãs.
Restrições: $5 \leq N \leq 100$

\OutputFile

Seu programa deve produzir uma única linha, contendo um único número inteiro, a idade de Camila.

% \Notes
% Preste atenção aos tipos das variáveis, pois $F(90) \approx 2^{61}$.

\Example

\begin{example}
\exmp{
6
9
7
}{
7
}%
\end{example}

\begin{example}
\exmp{
34
36
38
}{
36
}%
\end{example}

\begin{example}
\exmp{
22
25
22
}{
22
}%
\end{example}

\begin{example}
\exmp{
91
91
91
}{
91
}%
\end{example}

\textit{Questão da OBI2021, fase local.}

\end{problem}

\end{document}
